\scanpage{183}
\mspage{94}
% L01–09
Les points antipodiques à $s_0$ (resp. str\textsuperscript{t} antipodiques) sont donc les images de $s_0$ par les applications antipodiques (resp str. antipodiques). En un sens évident à expliciter, ces applications correspondent aussi biunivoquement aux ``types'' (à $G$-congruence près) de couples $(s_0,s)$ de points antipodiques (resp. strictement antipodiques).

% L09–14
On fera attention que la relation d'antipodisme entre pts n'est pas néc. une relation d'équivalence, tandis que la relation stricte l'est -- c'est la relation de $G$-équivalence envisagée au début.

% L15–21
Je m'intéresse maintenant au cas où $S$ est l'ensemble des sommets d'un 2-polyèdre \inserted{(combinatoire)} régulier \inserted{connexe}, \deleted{ayant $G$ \unsure{comme}} \inserted{et de son} groupe des automorphismes, \inserted{(noté maintenant $G_0$)} et à la variante orientée -- $G_0$ étant \inserted{alors} le groupe des autom orientés. Dans le cas où les ordres des sommets sont \unsure{finis|pairs},

% archive sheet number: 182
